\documentclass[12pt, a4paper]{article}

% ── Packages ──────────────────────────────────────────────
\usepackage[utf8]{inputenc}
\usepackage[T1]{fontenc}
\usepackage{geometry}
\usepackage{hyperref}
\usepackage{enumitem}
\usepackage{booktabs}
\usepackage{xcolor}
\usepackage{listings}
\usepackage{float}
\usepackage{amsmath}
\usepackage{array}

\geometry{margin=1in}

\hypersetup{
    colorlinks=true,
    linkcolor=blue,
    urlcolor=blue,
    citecolor=blue
}

\lstset{
    basicstyle=\ttfamily\small,
    backgroundcolor=\color{gray!10},
    frame=single,
    breaklines=true
}

% ── Title ─────────────────────────────────────────────────
\title{
    \textbf{Design Document} \\
    \Large Accelerated Mobile Video Processing System \\
}
\author{
    Salil Gujar (2025MCS2106) \\
    Nitish Kumar (2025MCS2100)
}
\date{}

% ══════════════════════════════════════════════════════════
\begin{document}

\maketitle


% ── 1. Introduction ──────────────────────────────────────
\section{Introduction}

\subsection{Project Overview}
This project aims to build a real-time Android camera application that applies per-pixel
filters — specifically a $5 \times 5$ Gaussian blur — to live video frames and displays
the result. The system must process live camera frames continuously and display the
transformed output in real time. It needs four processing modes: Baseline (plain C),
SIMD (ARM NEON), GPU (OpenGL ES compute shader), and Hybrid (CPU + GPU splitting
the frame). This is evaluated as a systems and performance engineering problem, not as
a graphics or UI exercise. All computation occurs entirely on the mobile device.

\subsection{Goals}
\begin{itemize}
    \item Implement a functional end-to-end camera pipeline using the Android Camera2 API,
          capable of continuous frame capture, processing, and display without manual triggering.
    \item Accelerate pixel processing using ARM NEON intrinsics and OpenGL ES Compute Shaders,
          achieving measurable throughput improvements over the scalar baseline.
    \item Support runtime switching between all four execution modes without restarting the
          application.
    \item Display real-time performance metrics — FPS, latency, active mode, and resource
          utilisation indicators — via an on-screen dashboard overlay.
    \item Verify functional correctness of all accelerated modes against the baseline output
          within acceptable numerical tolerance.
\end{itemize}

% ── 2. System Overview ───────────────────────────────────
\section{System Overview}

\subsection{High-Level Architecture}
The system operates as a continuous camera pipeline. Frames are captured at $1280 \times 720$
resolution in \texttt{YUV\_420\_888} format using the Camera2 API. Each frame is first
converted to RGBA, then dispatched to one of four selectable processing backends. The processed
frame is rendered to a \texttt{SurfaceView} for display. A transparent HUD overlay sits above
the rendered output and shows live performance statistics.

\subsection{Target Device}
Physical Android device, API 26+, ARM64.

\subsection{Processing Modes}
\begin{enumerate}
    \item \textbf{Baseline:} Plain C implementation using nested loops, called via JNI. This
          is the correctness reference and the performance floor against which all other modes
          are measured.
    \item \textbf{SIMD:} Vectorised implementation using ARM NEON intrinsics. The filter is
          rewritten to operate on 128-bit registers, processing 16 bytes per cycle.
    \item \textbf{GPU:} The filter is offloaded to an OpenGL ES 3.1 compute shader. Data goes
          to the GPU through an SSBO, the shader runs, and the result is read back.
    \item \textbf{Hybrid:} A split-frame approach — top half to NEON, bottom half to GPU —
          stitched before display.
\end{enumerate}

The output produced in all modes must be functionally equivalent within acceptable numerical
tolerance. Switching between execution modes must occur at runtime without restarting the
application.

% ── 3. Detailed Module Design ────────────────────────────
\section{Detailed Module Design}

\subsection{Camera Capture Module}
The rear-facing camera is opened via \texttt{CameraManager.openCamera()} with a
\texttt{CaptureRequest} configured for continuous preview at a fixed 30 fps
(\texttt{CONTROL\_AF\_MODE\_CONTINUOUS\_VIDEO}). Frames are delivered in \texttt{YUV\_420\_888}
format at $1280 \times 720$ through an \texttt{ImageReader}. An \texttt{OnImageAvailableListener}
processes each frame on a dedicated background \texttt{HandlerThread}. Camera and storage
permissions are declared in \texttt{AndroidManifest.xml}.

\subsection{YUV to RGBA Conversion}
In the JNI/C layer, the Y, U, and V planes and strides are extracted from the
\texttt{YUV\_420\_888} image and converted to RGBA per pixel using the standard BT.601 formula.
The result is returned to the Java/Kotlin rendering layer. This conversion is implemented in
native C to minimise overhead.

\subsection{Baseline Filter Module}
The baseline module implements the $5 \times 5$ Gaussian blur as nested loops in C, invoked
from Java via JNI. The kernel uses precomputed float weights normalised to 1.0. R, G, and B
channels are processed independently; the alpha channel passes through unchanged. Border pixels
use clamp-to-edge padding. This scalar implementation is intentionally simple — its sole purpose
is to be provably correct and to serve as the reference that all future accelerated paths are
verified against.

\subsection{SIMD Filter Module}
The SIMD module rewrites the Gaussian blur using ARM NEON intrinsics. The $5 \times 5$ kernel
is decomposed into two separable 1D passes — horizontal then vertical — reducing per-pixel
multiply operations from 25 to 10. The horizontal pass is implemented as follows:
\begin{itemize}
    \item \texttt{vld1q\_u8} loads 16-byte chunks
    \item \texttt{vextq\_u8} constructs shifted neighbour views for kernel access
    \item \texttt{vmull\_u8} and \texttt{vmlal\_u8} handle weighted accumulation
    \item \texttt{vshrn\_n\_u16} narrows back to \texttt{uint8x8\_t}
\end{itemize}
The vertical pass follows the same structure using a row pointer array. Pixel data buffers are
aligned to 16-byte boundaries to ensure optimal register utilisation. Border handling matches
the scalar reference exactly so that correctness comparison is unambiguous.

\subsection{GPU Filter Module}
The GPU module offloads the filter to an OpenGL ES 3.1 compute shader. An EGL context is
initialised on the processing thread (\texttt{EGL\_OPENGL\_ES3\_BIT}) with an offscreen
\texttt{EGL\_PBUFFER\_BIT} surface for headless GPU computation. The RGBA frame buffer is
uploaded to the GPU as an SSBO:

\begin{lstlisting}[language=C]
glGenBuffers(1, &ssbo);
glBindBuffer(GL_SHADER_STORAGE_BUFFER, ssbo);
glBufferData(GL_SHADER_STORAGE_BUFFER, size, data, GL_DYNAMIC_COPY);
\end{lstlisting}

The compute shader (\texttt{\#version 310 es}) assigns one invocation per pixel in $16 \times 16$
work groups. Each invocation reads a $5 \times 5$ neighbourhood from the input SSBO, applies
Gaussian weights, and writes its result to the output SSBO. Image borders are handled with
conditional clamping in the shader. The dispatch is:

\begin{lstlisting}[language=C]
glDispatchCompute(width / 16, height / 16, 1);
glMemoryBarrier(GL_SHADER_STORAGE_BARRIER_BIT);
\end{lstlisting}

Results are read back to CPU via \texttt{glGetBufferSubData()}. The compute shader is a
self-contained program and is developed independently of the CPU paths.

\subsection{Hybrid Filter Module}
The hybrid module manages workload distribution between the NEON and GPU paths within a single
frame. The $1280 \times 720$ frame is split horizontally at a configurable split row (default:
row 360). The top half is sent to the NEON filter on a CPU \texttt{pthread}; the bottom half
is dispatched to the GPU shader concurrently. Completion is synchronised via
\texttt{pthread\_join()} for the NEON thread and \texttt{glFinish()} on the GL thread. The
two output halves are stitched with \texttt{memcpy()}. The split ratio is configurable and
displayed on the dashboard.

Rather than a fixed 50/50 split, the split row is computed dynamically based on measured
throughput:
\[
    \text{split\_row} = \frac{\text{NEON\_throughput}}{\text{NEON\_throughput} + \text{GPU\_throughput}} \times \text{height}
\]
This balances both execution paths to finish at approximately the same time, eliminating the
idle wait before stitching.

\subsection{Dashboard / HUD Module}
The dashboard is a mandatory system component and is treated with the same rigour as the
processing backends. It is rendered as a semi-transparent overlay above the live camera output
and updates every rendered frame. The dashboard displays:
\begin{itemize}
    \item Currently active execution mode (\texttt{BASELINE / SIMD / GPU / HYBRID})
    \item Live frame rate (FPS) — rolling average over last 30 frames
    \item Per-frame processing latency in milliseconds
    \item End-to-end latency from \texttt{ImageReader} callback to \texttt{unlockCanvasAndPost()}
    \item SIMD active indicator — green (active) / red (inactive)
    \item GPU execution activity indicator — active when a shader is dispatched
    \item Thermal proxy — CPU temperature in \textdegree C, polled from
          \texttt{/sys/class/thermal/thermal\_zone*/temp}
    \item Speedup ratio vs.\ baseline for all non-baseline modes (e.g., $3.2\times$)
    \item Jitter events — count of frame time spikes exceeding $2\times$ the rolling median
          in the last 60 seconds
    \item Frame split ratio in Hybrid mode (e.g., \texttt{360 rows CPU / 360 rows GPU})
    \item PASS/FAIL correctness flags for SIMD and GPU modes
\end{itemize}

The final dashboard layout is as follows:
\begin{lstlisting}
+---------------------------------------------+
|  MODE: [HYBRID]         FPS: 42.3           |
|  Frame latency:  18.4 ms                    |
|  End-to-end:     22.1 ms                    |
+---------------------------------------------+
|  SIMD  [################    ]  ACTIVE       |
|  GPU   [########            ]  ACTIVE       |
|  Split: 360 rows CPU / 360 rows GPU         |
+---------------------------------------------+
|  Baseline: 84.2 ms   Speedup: 4.6x          |
|  Thermal: 38 C        Jitter events: 0      |
+---------------------------------------------+
\end{lstlisting}

Mode label colours: grey $\rightarrow$ Baseline, blue $\rightarrow$ SIMD, green $\rightarrow$
GPU, yellow $\rightarrow$ Hybrid.

% ── 4. Data Flow ─────────────────────────────────────────
\section{Data Flow and Frame Lifecycle}

Each frame follows the same pipeline regardless of the active execution mode:
\begin{enumerate}
    \item \textbf{Capture:} Camera2 API delivers a \texttt{YUV\_420\_888} frame via
          \texttt{ImageReader} callback.
    \item \textbf{Conversion:} Frame is converted from YUV to a packed RGBA buffer in native
          C via JNI.
    \item \textbf{Dispatch:} The RGBA buffer is handed to the currently active backend. In
          Hybrid mode, the buffer is split and dispatched to two backends simultaneously.
    \item \textbf{Processing:} The selected backend applies the $5 \times 5$ Gaussian blur and
          writes the result into an output buffer.
    \item \textbf{Stitch (Hybrid only):} The NEON and GPU output halves are merged into a
          single output buffer via \texttt{memcpy()}.
    \item \textbf{Render:} The output buffer is drawn to the \texttt{SurfaceView} via
          \texttt{Bitmap.copyPixelsFromBuffer()} and \texttt{Canvas}, and the dashboard
          overlay is updated with the latest metrics.
\end{enumerate}

End-to-end latency is measured at four points in this pipeline: (1) \texttt{ImageReader}
callback entry, (2) YUV$\rightarrow$RGBA conversion complete, (3) filter processing complete,
and (4) \texttt{unlockCanvasAndPost()} complete. End-to-end latency is defined as point (4)
minus point (1) for every frame.

% ── 5. Correctness Verification ──────────────────────────
\section{Correctness Verification}

Both accelerated paths are verified against the scalar baseline before being made available
at runtime. The demo does not proceed until both SIMD and GPU correctness flags show PASS.

\textbf{SIMD:} Both scalar and NEON paths are run on the same static test frame. Outputs are
compared byte-by-byte in JNI using \texttt{memcmp()}, with a tolerance of $\pm 1$ per channel
to account for rounding differences in integer NEON arithmetic. SIMD mode is only activated
once verification passes. This gate remains permanent throughout the project.

\textbf{GPU:} The scalar baseline and GPU path are run on the same static frame and compared
in JNI. A tolerance of $\pm 2$ per channel is used to account for floating-point arithmetic
in the shader. GPU correctness status is shown on the dashboard.

% ── 6. Memory and Buffer Design ──────────────────────────
\section{Memory and Buffer Design}

To avoid garbage collection pauses in the hot path, three allocations are pre-allocated at
session start and reused across frames:
\begin{itemize}
    \item \textbf{Bitmap:} A single \texttt{Bitmap} is allocated once and updated in-place
          per frame via \texttt{copyPixelsFromBuffer()}, replacing per-frame
          \texttt{Bitmap.createBitmap()} calls.
    \item \textbf{SSBO:} The GPU Shader Storage Buffer Object is pre-allocated at session start
          and updated per frame with \texttt{glBufferSubData()} instead of
          \texttt{glBufferData()}.
    \item \textbf{NEON buffers:} Intermediate row buffers for the separable filter passes are
          pre-allocated; \texttt{malloc} is removed from the filter inner loop.
\end{itemize}





\end{document}